\section{Methodology}
\label{sec:method}

\subsection{Problem Formulation and Architecture Overview}
\label{sec:method:problem-overview}

We consider binary prediction over a relational database. At row level, the
complete database is represented by the directed multigraph
\begin{equation}
  \mathcal G^{\mathrm{DB}}
  =
  \bigl(\mathcal V^{\mathrm{DB}},\mathcal E^{\mathrm{FK}}\bigr),
\label{eq:method:database-graph}
\end{equation}
where \(\mathcal V^{\mathrm{DB}}\) is the set of database rows and
\(\mathcal E^{\mathrm{FK}}\) is an edge-index set induced by foreign-key
(FK) and primary-key references. Every \(e\in\mathcal E^{\mathrm{FK}}\) is
directed from a child row \(\rgsrc(e)\) to the referenced parent row
\(\rgdst(e)\). We treat edges as indexed objects rather than endpoint pairs,
so parallel FK links between the same two rows remain distinct.

Let \(\mathfrak T\) denote the collection of binary prediction tasks. A
query instance \(g\) belongs to task \(\kappa(g)\in\mathfrak T\) and is
specified by a seed row \(r_g^\star\in\mathcal V^{\mathrm{DB}}\), a target
column \(c_g^\star\), and a binary label \(y_g\in\{0,1\}\). We denote the
target cell located at \((r_g^\star,c_g^\star)\) by \(n_g^\star\). Its
\emph{value} is masked and excluded from all cell and row encodings. Its
structural metadata---the target column and semantic type, the seed table,
and the seed timestamp---remain available because they specify the query
rather than reveal its answer.
The task identifier \(\kappa(g)\) determines this target schema information
but does not select a task-specific parameter set. We present the binary
formulation first; the not-yet-finalized regression extension is reserved
for Section~\ref{sec:method:regression-placeholder}.

For each query, the model performs four stages. First, it samples a bounded
relational context centered at \(r_g^\star\). Second, it constructs two
query-specific graphs: a candidate-augmented row graph
\(\mathcal G_g^{\graphmain}\) and an auxiliary relation graph
\(\mathcal G_g^{\graphrel}\). Third, \(P\) relation-graph layers encode how
the sampled FK relation types interact and produce a condition vector for
every retained FK edge; \(L\) row-graph layers then propagate typed messages
among row and candidate nodes under those conditions. Finally, a shared
readout compares the final seed representation with each candidate
representation and converts their score difference into a binary
probability. The base model uses the same row graph, initialization,
message-passing operator, and readout, but sets every FK condition to zero.

We use \(d\) for the common hidden width, \(d_{\mathrm{text}}\) for the
dimension of precomputed textual features, \(F\) for the number of temporal
frequencies, \(P\) for the number of relation-graph layers, and \(L\) for
the number of row-graph layers. The classification configuration is
\begin{equation}
  d=256,
  \qquad
  d_{\mathrm{text}}=384,
  \qquad
  F=16,
  \qquad
  P=2,
  \qquad
  L=6.
\label{eq:method:model-dimensions}
\end{equation}

Throughout this section, \(g\) indexes a query, \(r\) a database row,
\(n\) a cell, \(k\) a candidate, and \(i,j\) generic row-graph nodes.
Row-graph edges and layers are indexed by \(e\) and \(\ell\), whereas
relation-graph nodes, edges, and layers are indexed by \(\rho\),
\(\epsilon\), and \(p\). Bold lowercase symbols denote vectors, bold
uppercase symbols denote matrices, calligraphic symbols denote sets or
graphs, and unbold symbols denote scalars, indices, or discrete objects.
All learnable parameters below are shared across tasks and query instances;
therefore, no trainable parameter carries \(g\) or \(\kappa(g)\) as an
index. Parameters are distinct across semantic types, edge types, and
layers whenever such an index is shown.

\subsection{Seed-Centered Relational Context Sampling}
\label{sec:method:context-sampling}

Applying the model directly to \(\mathcal G^{\mathrm{DB}}\) is generally
intractable. For each query \(g\), we therefore run a bounded,
seed-centered, priority-based hybrid traversal beginning at
\(r_g^\star\). The traversal follows both foreign-to-primary
(\(\mathsf{f2p}\)) and primary-to-foreign (\(\mathsf{p2f}\)) links and
maintains a visited-row set so that a row is serialized at most once.
Foreign-to-primary expansions use a prioritized last-in--first-out frontier
and can therefore follow a chain of referenced parents depth-first. Once
that frontier is empty, primary-to-foreign neighbors are taken from the
shallowest nonempty traversal depth and randomized within that depth.

Let the resulting non-padding cell sequence be
\begin{equation}
  \mathcal S_g
  =
  \operatorname{Sample}
  \left(
    \mathcal G^{\mathrm{DB}},
    r_g^\star;
    S_{\max},B_{\max}
  \right),
  \qquad
  |\mathcal S_g|\le S_{\max},
\label{eq:method:context-sample}
\end{equation}
where \(S_{\max}=2048\) is a \emph{cell} budget, not a row budget. Cells
are appended row by row until the budget is exhausted, so the last retained
row may be represented by only a prefix of its cells. Shorter sequences are
padded only after the non-padding prefix has been constructed. The target
cell can occupy a position in this serialized prefix, but its mask prevents
it from entering \(\mathcal C_g^{\mathrm{in}}\) or any subsequent pooling
operation.

To prevent a high-degree primary row from exhausting the context,
one reverse-FK expansion retains at most \(B_{\max}=256\) randomly selected
database children. For forecasting queries, a reverse-FK neighbor is
discarded when both its timestamp and the seed timestamp are present and
the neighbor occurs after the seed. This causal filter is applied during
primary-to-foreign expansion rather than globally to every sampled row.
Task-defined leakage columns are also removed from the seed and from rows
whose stored optional timestamp equals the seed timestamp; the
implementation's equality rule includes the case in which both timestamps
are missing.

The sampler additionally exposes local-recency, peripheral-cell-cap, and
self-label-dropout controls. They change which context cells remain visible
but do not change the graph neural architecture defined below. The default
hybrid traversal described above uses no additional local-recency quota;
any nondefault sampling or label-dropout setting must therefore be reported
as part of the experimental protocol rather than treated as a new model
component.

The cells actually supplied to the encoder are
\begin{equation}
  \mathcal C_g^{\mathrm{in}}
  =
  \left\{
    n\in\mathcal S_g
    \,\middle|\,
    n\text{ is not masked}
  \right\}.
\label{eq:method:input-cell-set}
\end{equation}
Let \(r(n)\) and \(c(n)\) denote the row and column containing cell \(n\).
The retained row set and the observed cells belonging to one retained row
are, respectively,
\begin{equation}
\begin{aligned}
  \mathcal R_g
  &=
  \{r(n):n\in\mathcal C_g^{\mathrm{in}}\},\\
  \mathcal O_{g,r}
  &=
  \{n\in\mathcal C_g^{\mathrm{in}}:r(n)=r\},
  \qquad r\in\mathcal R_g.
\end{aligned}
\label{eq:method:sampled-rows-and-cells}
\end{equation}
By construction, \(\mathcal O_{g,r}\neq\varnothing\) for every represented
row and
\begin{equation}
  n_g^\star
  \notin
  \bigcup_{r\in\mathcal R_g}\mathcal O_{g,r}.
\label{eq:method:target-exclusion}
\end{equation}
Thus, the target value cannot enter the model through either cell encoding
or row pooling. The complete database graph in
Eq.~\eqref{eq:method:database-graph} is used only for context selection;
neural message passing operates on the two finite query-specific graphs
defined next.

\subsection{Query-Specific Graph Construction}
\label{sec:method:graph-construction}

\subsubsection{Candidate-Augmented Row Graph}
\label{sec:method:row-graph}

For query \(g\), the main prediction graph is the directed, typed
multigraph
\begin{equation}
  \mathcal G_g^{\graphmain}
  =
  \bigl(
    \mathcal V_g^{\graphmain},
    \mathcal E_g^{\graphmain},
    \tau_g,
    \alpha_g;
    r_g^\star
  \bigr).
\label{eq:method:main-graph}
\end{equation}
The semicolon emphasizes that \(r_g^\star\) is a distinguished query node
rather than a separate node type. The node set is the disjoint union
\begin{equation}
\begin{aligned}
  \mathcal V_g^{\graphmain}
  &=
  \mathcal V_g^{\mathrm{row}}
  \mathbin{\dot\cup}
  \mathcal V_g^{\mathrm{cand}},\\
  \mathcal V_g^{\mathrm{row}}
  &=
  \mathcal R_g,
  &
  \mathcal V_g^{\mathrm{cand}}
  &=
  \{v_{g,0}^{\mathrm{cand}},v_{g,1}^{\mathrm{cand}}\}.
\end{aligned}
\label{eq:method:main-node-set}
\end{equation}
The candidate nodes \(v_{g,0}^{\mathrm{cand}}\) and
\(v_{g,1}^{\mathrm{cand}}\) represent the hypotheses \textsc{False} and
\textsc{True}, respectively. The target cell itself is not a graph node:
\(n_g^\star\notin\mathcal V_g^{\graphmain}\), whereas its containing seed
row satisfies \(r_g^\star\in\mathcal V_g^{\mathrm{row}}\).

The edge-index set and available edge types are
\begin{equation}
\begin{aligned}
  \mathcal E_g^{\graphmain}
  &=
  \mathcal E_g^{\mathsf{f2p}}
  \mathbin{\dot\cup}
  \mathcal E_g^{\mathsf{p2f}}
  \mathbin{\dot\cup}
  \mathcal E_g^{\mathsf{label}},\\
  \mathcal T_{\graphmain}
  &=
  \{\mathsf{f2p},\mathsf{p2f},\mathsf{label}\}.
\end{aligned}
\label{eq:method:main-edge-set}
\end{equation}
The maps
\begin{equation}
  \tau_g:\mathcal E_g^{\graphmain}\to\mathcal T_{\graphmain},
  \qquad
  \alpha_g:\mathcal E_g^{\graphmain}\to\RR_{\ge0}
\label{eq:method:main-edge-maps}
\end{equation}
assign each edge \(e\) its type \(\tau_e\coloneqq\tau_g(e)\) and fixed
scalar weight \(\alpha_e\coloneqq\alpha_g(e)\). These weights are part of
the graph input, not trainable parameters. For binary classification,
\(\alpha_e=1\) for all three edge types.

\subsubsection{Foreign-Key and Observed-Label Edges}
\label{sec:method:row-graph-edges}

The foreign-to-primary edges are the stored database FK edges whose two
endpoints both survive context sampling:
\begin{equation}
  \mathcal E_g^{\mathsf{f2p}}
  =
  \left\{
    e\in\mathcal E^{\mathrm{FK}}
    \,\middle|\,
    \rgsrc(e),\rgdst(e)\in\mathcal R_g
  \right\}.
\label{eq:method:f2p-edges}
\end{equation}
Here \(\rgsrc(e)\) is a sampled child row and \(\rgdst(e)\) is its sampled
parent. The fixed-width stored representation requires at most five
foreign-to-primary neighbors per row. For every retained forward edge
\(e\), we create a distinct reverse edge \(\bar e\):
\begin{equation}
  \mathcal E_g^{\mathsf{p2f}}
  =
  \left\{
    \bar e
    \,\middle|\,
    \begin{array}{l}
      e\in\mathcal E_g^{\mathsf{f2p}},\\[-1mm]
      \rgsrc(\bar e)=\rgdst(e),\quad
      \rgdst(\bar e)=\rgsrc(e)
    \end{array}
  \right\}.
\label{eq:method:p2f-edges}
\end{equation}
Consequently, \(\tau_e=\mathsf{f2p}\),
\(\tau_{\bar e}=\mathsf{p2f}\), and
\(\alpha_e=\alpha_{\bar e}=1\). The two directions preserve the same
underlying FK identity but use different message transformations.

Observed-label edges inject visible, target-aligned supervision from the
sampled context into the output candidates. Define
\begin{equation}
  \mathcal R_g^{\mathrm{lab}}
  =
  \left\{
    r\in\mathcal R_g\setminus\{r_g^\star\}
    \,\middle|\,
    \begin{array}{l}
      \tblname(r)=\tblname(r_g^\star),\\[-1mm]
      \exists n\in\mathcal O_{g,r}:c(n)=c_g^\star
    \end{array}
  \right\}.
\label{eq:method:observed-label-rows}
\end{equation}
For \(r\in\mathcal R_g^{\mathrm{lab}}\), let
\(y_{g,r}\in\{0,1\}\) denote the visible Boolean value in column
\(c_g^\star\). This contextual label is distinct from the unknown query
label \(y_g\). We add one directed edge to the matching candidate:
\begin{equation}
  \mathcal E_g^{\mathsf{label}}
  =
  \left\{
    e_{g,r}^{\mathrm{label}}
    =
    \bigl(r,v_{g,y_{g,r}}^{\mathrm{cand}}\bigr)
    \,\middle|\,
    r\in\mathcal R_g^{\mathrm{lab}}
  \right\},
\label{eq:method:label-edges}
\end{equation}
where
\(\tau_{e_{g,r}^{\mathrm{label}}}=\mathsf{label}\) and
\(\alpha_{e_{g,r}^{\mathrm{label}}}=1\). No candidate-to-row reverse edge
is created, and candidate nodes have no outgoing edges in the current
construction. Because the seed is excluded from
\(\mathcal R_g^{\mathrm{lab}}\), its target cannot create a label edge.
Thus, label edges act as in-graph demonstrations without exposing the
answer to the current query. Disabling observed-label edges for an ablation
sets \(\mathcal E_g^{\mathsf{label}}=\varnothing\) without changing the FK
edges or the propagation equations.

\subsubsection{Auxiliary Relation Graph}
\label{sec:method:relation-graph}

The generic row-graph types \(\mathsf{f2p}\) and \(\mathsf{p2f}\) identify
direction but not the schema-level role of an FK link. We therefore define
the auxiliary directed, typed multigraph
\begin{equation}
  \mathcal G_g^{\graphrel}
  =
  \bigl(
    \mathcal V_g^{\graphrel},
    \mathcal E_g^{\graphrel},
    \eta_g;
    \mathcal Q_g^{\graphrel}
  \bigr).
\label{eq:method:relation-graph}
\end{equation}
For each \(e\in\mathcal E_g^{\mathsf{f2p}}\), its ordered table-pair
relation is
\begin{equation}
  \chi_g(e)
  =
  \bigl(
    \tblname(\rgsrc(e)),
    \tblname(\rgdst(e))
  \bigr),
\label{eq:method:fk-to-relation}
\end{equation}
where the first component is the child table and the second is the parent
table. The relation-node set contains the unique table pairs present in the
sampled row graph:
\begin{equation}
  \mathcal V_g^{\graphrel}
  =
  \{\chi_g(e):e\in\mathcal E_g^{\mathsf{f2p}}\}.
\label{eq:method:relation-node-set}
\end{equation}
Multiple row-level links with the same ordered table pair share one
query-local relation node \(\rho\); relation nodes from different queries
\(g\ne g'\) are never merged. The current construction deliberately
collapses multiple FK columns connecting the same ordered pair of tables.
For the reverse copy of a retained FK link, we define
\begin{equation}
  \chi_g(\bar e)\coloneqq\chi_g(e),
  \qquad e\in\mathcal E_g^{\mathsf{f2p}},
\label{eq:method:reverse-fk-relation}
\end{equation}
so the forward and reverse directions map to the same relation node.

For \(\rho\in\mathcal V_g^{\graphrel}\), let
\(\operatorname{child}(\rho)\) and \(\operatorname{parent}(\rho)\) denote
the two components of its ordered table pair. For distinct relation nodes
\(\rho,\rho'\), four predicates define directed structural edge types:
\begin{equation}
\begin{aligned}
  \mathsf{match}_{\mathsf{h2h}}(\rho,\rho')
  &:\ \operatorname{child}(\rho)=\operatorname{child}(\rho'),
  &
  \mathsf{match}_{\mathsf{t2t}}(\rho,\rho')
  &:\ \operatorname{parent}(\rho)=\operatorname{parent}(\rho'),\\
  \mathsf{match}_{\mathsf{h2t}}(\rho,\rho')
  &:\ \operatorname{child}(\rho)=\operatorname{parent}(\rho'),
  &
  \mathsf{match}_{\mathsf{t2h}}(\rho,\rho')
  &:\ \operatorname{parent}(\rho)=\operatorname{child}(\rho').
\end{aligned}
\label{eq:method:relation-predicates}
\end{equation}
Here \emph{head} denotes the child table and \emph{tail} the parent table.
With
\begin{equation}
  \mathcal T_{\graphrel}
  =
  \{\mathsf{h2h},\mathsf{t2t},\mathsf{h2t},\mathsf{t2h}\},
\label{eq:method:relation-edge-types}
\end{equation}
the complete relation-edge index set is
\begin{equation}
  \mathcal E_g^{\graphrel}
  =
  \left\{
    (\rho,\rho',\eta)
    \,\middle|\,
    \rho\ne\rho',\ 
    \eta\in\mathcal T_{\graphrel},\ 
    \mathsf{match}_\eta(\rho,\rho')
  \right\}.
\label{eq:method:relation-edge-set}
\end{equation}
The map
\(\eta_g:\mathcal E_g^{\graphrel}\to\mathcal T_{\graphrel}\) assigns
each relation edge \(\epsilon\) its type
\(\eta_\epsilon\coloneqq\eta_g(\epsilon)\). All relation-edge weights are
one. If multiple predicates hold for the same ordered endpoint pair, the
implementation retains parallel edges carrying different type tags.
When the sampled row graph contains no FK edge,
\(\mathcal V_g^{\graphrel}=\mathcal E_g^{\graphrel}=\varnothing\), and the
relation branch is skipped.

\subsubsection{Query-Relation Marking}
\label{sec:method:query-relation-marking}

The relation graph becomes query specific through the distinguished set of
relations in which the seed appears as the FK child/source:
\begin{equation}
  \mathcal Q_g^{\graphrel}
  =
  \left\{
    \rho\in\mathcal V_g^{\graphrel}
    \,\middle|\,
    \begin{array}{l}
      \exists e\in\mathcal E_g^{\mathsf{f2p}}:\\[-1mm]
      \rgsrc(e)=r_g^\star\ \text{and}\ \chi_g(e)=\rho
    \end{array}
  \right\}.
\label{eq:method:query-relation-set}
\end{equation}
A relation is therefore not marked merely because the seed is its parent.
One seed can mark multiple relations, and no relation is marked when the
seed has no retained outgoing FK edge.
The set \(\mathcal Q_g^{\graphrel}\) supplies the boundary signal used to
initialize relation-node states in
Section~\ref{sec:method:candidate-relation-initialization}; it is the only
query-boundary marker provided to the auxiliary relation graph.

\subsection{Initial Node Representations}
\label{sec:method:node-initialization}

\subsubsection{Type-Aware Cell Encoding}
\label{sec:method:cell-encoding}

Let \(n\in\mathcal O_{g,r}\) be an observed input cell in row \(r\), and
let
\begin{equation}
  \sigma_{r,n}
  \in
  \Sigma_{\mathrm{sem}}
  =
  \{\mathsf{num},\mathsf{text},\mathsf{datetime},\mathsf{bool}\}
\label{eq:method:semantic-types}
\end{equation}
denote its semantic type. Each cell supplies a preprocessed value feature
\(\mathbf u_{r,n}^{\mathrm{val}}\in\RR^{d_{\sigma_{r,n}}}\) and a
precomputed column-name feature
\(\mathbf u_{r,n}^{\mathrm{col}}\in\RR^{d_{\mathrm{text}}}\). Text values
and column names use \(d_{\mathrm{text}}=384\)-dimensional text features;
numeric, datetime, and Boolean values are represented by scalar inputs.
Thus,
\begin{equation}
  d_\sigma
  =
  \begin{cases}
    d_{\mathrm{text}}, & \sigma=\mathsf{text},\\
    1, & \sigma\in\{\mathsf{num},\mathsf{datetime},\mathsf{bool}\}.
  \end{cases}
\label{eq:method:value-input-dimensions}
\end{equation}

A type-specific affine map encodes the value, while a single shared affine
map encodes the column name. Their normalized outputs are added:
\begin{equation}
\begin{aligned}
  \mathbf x_{r,n}
  ={}&
  \rmsnorm_{\mathrm{col}}
  \left(
    \mathbf W_{\mathrm{col}}\mathbf u_{r,n}^{\mathrm{col}}
    +\mathbf b_{\mathrm{col}}
  \right)\\
  &+
  \rmsnorm_{\sigma_{r,n}}
  \left(
    \mathbf W_{\mathrm{val},\sigma_{r,n}}
    \mathbf u_{r,n}^{\mathrm{val}}
    +\mathbf b_{\mathrm{val},\sigma_{r,n}}
  \right)
  \in\RR^d.
\end{aligned}
\label{eq:method:cell-encoding}
\end{equation}
The trainable affine parameters have domains
\begin{equation}
\begin{aligned}
  \mathbf W_{\mathrm{col}}
  &\in\RR^{d\times d_{\mathrm{text}}},
  &
  \mathbf b_{\mathrm{col}}
  &\in\RR^d,\\
  \mathbf W_{\mathrm{val},\sigma}
  &\in\RR^{d\times d_\sigma},
  &
  \mathbf b_{\mathrm{val},\sigma}
  &\in\RR^d,
  \qquad \sigma\in\Sigma_{\mathrm{sem}}.
\end{aligned}
\label{eq:method:cell-parameter-domains}
\end{equation}
Value encoders are distinct across semantic types; the column-name encoder
is shared across all columns, rows, tasks, and queries.

For \(\boldsymbol\zeta\in\RR^d\), every RMSNorm module indexed by
\(\nu\in\{\mathrm{col}\}\cup\Sigma_{\mathrm{sem}}\) is
\begin{equation}
  \rmsnorm_\nu(\boldsymbol\zeta)
  =
  \boldsymbol\gamma_\nu\odot
  \frac{\boldsymbol\zeta}{
    \sqrt{d^{-1}\|\boldsymbol\zeta\|_2^2+\varepsilon_{\mathrm{rms}}}
  },
  \qquad
  \boldsymbol\gamma_\nu\in\RR^d,
\label{eq:method:rmsnorm}
\end{equation}
where \(\odot\) is element-wise multiplication and
\(\varepsilon_{\mathrm{rms}}>0\) is a numerical stabilizer. Each module has
its own learned scale \(\boldsymbol\gamma_\nu\). Separate value encoders and
normalizers handle heterogeneous value modalities; the column-name term
distinguishes fields whose values look similar but whose semantics differ.
Because \(n_g^\star\notin\mathcal C_g^{\mathrm{in}}\), neither the target
value feature nor a target-cell representation contributes to the seed row.

\subsubsection{Permutation-Invariant Cell-to-Row Pooling}
\label{sec:method:row-pooling}

For each \(r\in\mathcal R_g\), mean pooling summarizes the complete set of
observed cells:
\begin{equation}
  \overline{\mathbf x}_{g,r}
  =
  \frac{1}{|\mathcal O_{g,r}|}
  \sum_{n\in\mathcal O_{g,r}}\mathbf x_{r,n}
  \in\RR^d.
\label{eq:method:mean-pooling}
\end{equation}
Coordinate-wise max pooling retains the strongest response in each hidden
dimension:
\begin{equation}
  [\mathbf x_{g,r}^{\max}]_q
  =
  \max_{n\in\mathcal O_{g,r}}[\mathbf x_{r,n}]_q,
  \qquad q=1,\ldots,d.
\label{eq:method:max-pooling}
\end{equation}
Their concatenation is projected back to the common hidden width:
\begin{equation}
  \overline{\mathbf h}_{g,r}
  =
  \mathbf W_{\mathrm{c2r}}
  [\overline{\mathbf x}_{g,r}\vcat\mathbf x_{g,r}^{\max}]
  +\mathbf b_{\mathrm{c2r}}
  \in\RR^d,
\label{eq:method:cell-to-row}
\end{equation}
with
\begin{equation}
  \mathbf W_{\mathrm{c2r}}\in\RR^{d\times2d},
  \qquad
  \mathbf b_{\mathrm{c2r}}\in\RR^d.
\label{eq:method:c2r-parameter-domains}
\end{equation}
Mean pooling captures overall row content, whereas max pooling preserves
salient cell-level signals that could be diluted by averaging. Since both
operate on the set \(\mathcal O_{g,r}\), the resulting row representation
is invariant to the order in which cells were serialized. In the default
model, Eq.~\eqref{eq:method:cell-to-row} is a single affine projection with
no intervening ReLU or normalization; attention pooling is an optional
ablation rather than part of the mean--max architecture specified here.

\subsubsection{Relative-Time and Seed Encoding}
\label{sec:method:time-seed-encoding}

Let \(T_r\) be the optional timestamp of row \(r\), stored in seconds. We
express time relative to the seed in days:
\begin{equation}
  \Delta t_{g,r}
  =
  \frac{T_{r_g^\star}-T_r}{86400}.
\label{eq:method:relative-time}
\end{equation}
The division by \(86400\) converts seconds to days, matching the units of
the frequency bank. If either required timestamp is unavailable, the
implementation uses \(\Delta t_{g,r}=0\).

For \(f=0,\ldots,F-1\), define fixed wavelengths and angular frequencies
by
\begin{equation}
  \lambda_f
  =
  10^{4.56f/(F-1)}\ \text{days},
  \qquad
  \omega_f
  =
  \frac{2\pi}{\lambda_f}.
\label{eq:method:temporal-frequencies}
\end{equation}
Writing
\(\boldsymbol\omega=[\omega_0,\ldots,\omega_{F-1}]^\top\in\RR^F\), the
multi-scale temporal feature is
\begin{equation}
  \boldsymbol\phi_{\mathrm{time}}(\Delta t_{g,r})
  =
  [
    \cos(\Delta t_{g,r}\boldsymbol\omega)
    \vcat
    \sin(\Delta t_{g,r}\boldsymbol\omega)
  ]
  \in\RR^{2F},
\label{eq:method:time-feature}
\end{equation}
where the trigonometric functions act element-wise. The frequency vector is
a fixed buffer, not a learned parameter. With \(F=16\), the wavelengths are
logarithmically distributed from one day to approximately one hundred
years; short wavelengths resolve nearby times, while long wavelengths vary
slowly enough to distinguish long horizons.

The initial representation of row node \(r\) is
\begin{equation}
  \mathbf h_{g,r}^{(0)}
  =
  \overline{\mathbf h}_{g,r}
  +\mathbf W_{\mathrm{time}}
   \boldsymbol\phi_{\mathrm{time}}(\Delta t_{g,r})
  +\mathbf b_{\mathrm{time}}
  +\Indic[r=r_g^\star]\mathbf e_{\mathrm{seed}}
  \in\RR^d,
\label{eq:method:row-initialization}
\end{equation}
where
\begin{equation}
  \mathbf W_{\mathrm{time}}\in\RR^{d\times2F},
  \qquad
  \mathbf b_{\mathrm{time}}\in\RR^d,
  \qquad
  \mathbf e_{\mathrm{seed}}\in\RR^d.
\label{eq:method:time-seed-parameter-domains}
\end{equation}
The same learned temporal projection is shared by every row, task, and
query. The learned marker \(\mathbf e_{\mathrm{seed}}\) is added only to the
current seed and makes the prediction boundary explicitly visible to the
row GNN. Notice that a missing timestamp sets \(\Delta t_{g,r}=0\), not the
entire temporal contribution to zero: the row still receives
\(\mathbf W_{\mathrm{time}}\boldsymbol\phi_{\mathrm{time}}(0)
+\mathbf b_{\mathrm{time}}\).

\subsubsection{Candidate- and Relation-Node Initialization}
\label{sec:method:candidate-relation-initialization}

The two classification candidates are represented by the fixed scalar
values \(0\) and \(1\). To distinguish a candidate index from a generic
graph-node index, define the shorthand
\begin{equation}
  \mathbf h_{g,k}^{\mathrm{cand},(\ell)}
  \coloneqq
  \mathbf h_{g,v_{g,k}^{\mathrm{cand}}}^{(\ell)},
  \qquad k\in\{0,1\}.
\label{eq:method:candidate-state-shorthand}
\end{equation}
The initial candidate state reuses exactly the Boolean value affine map and
RMSNorm from Eq.~\eqref{eq:method:cell-encoding}:
\begin{equation}
  \mathbf h_{g,k}^{\mathrm{cand},(0)}
  =
  \rmsnorm_{\mathsf{bool}}
  \left(
    \mathbf W_{\mathrm{val},\mathsf{bool}}[k]
    +\mathbf b_{\mathrm{val},\mathsf{bool}}
  \right)
  +\mathbf e_{\mathrm{cand}}
  \in\RR^d,
  \qquad k\in\{0,1\},
\label{eq:method:candidate-initialization}
\end{equation}
where \([k]\in\RR^1\),
\(\mathbf W_{\mathrm{val},\mathsf{bool}}\in\RR^{d\times1}\),
\(\mathbf b_{\mathrm{val},\mathsf{bool}}\in\RR^d\), and
\(\mathbf e_{\mathrm{cand}}\in\RR^d\). Therefore, False and True are not
two independent learned embedding vectors. They are canonical Boolean
values passed through a shared encoder and augmented by the same learned
candidate-type marker. Candidate initialization contains neither a target
column embedding nor a task-specific embedding. Before graph propagation,
the same candidate \(k\) has the same initial representation in every
query; its representation becomes query specific only through subsequent
message passing.

For relation node \(\rho\in\mathcal V_g^{\graphrel}\), the initial state is
\begin{equation}
  \mathbf z_{g,\rho}^{(0)}
  =
  \Indic[\rho\in\mathcal Q_g^{\graphrel}]
  \mathbf e_{\mathrm{relquery}}
  \in\RR^d,
\label{eq:method:relation-initialization}
\end{equation}
where \(\mathbf e_{\mathrm{relquery}}\in\RR^d\) is a learned marker shared
across relation nodes, tasks, and queries. Query-relation nodes receive this
marker; all other relation nodes start at zero. No trainable schema-ID or
table-pair embedding is assigned to a relation node, so relation identity is
expressed through graph structure and the query marker rather than a fixed
lookup table.

The marker \(\mathbf e_{\mathrm{seed}}\) is shared by all seed rows,
\(\mathbf e_{\mathrm{cand}}\) by both candidates, and
\(\mathbf e_{\mathrm{relquery}}\) by all marked relation nodes. In the
implementation, the seed and candidate markers are initialized to zero,
whereas the relation-query marker is initialized with standard deviation
\(0.02\); all are trainable thereafter. The main-graph and relation-graph
propagation networks use separate parameters. Within either network,
parameters carrying an edge-type or layer index are not shared across
different values of that index.

\subsection{Relation-Conditioned Message Passing}
\label{sec:method:message-passing}

\subsubsection{Relation-Graph Propagation}
\label{sec:method:relation-propagation}

The auxiliary relation graph is processed completely before any row-graph
layer is evaluated. Let \(\epsilon:\rho'\to\rho\) be a directed relation
edge with type
\(\eta_\epsilon\equiv\eta_g(\epsilon)\in\mathcal T_{\graphrel}\). At
relation layer \(p\in\{1,\ldots,P\}\), its message is
\begin{equation}
  \widetilde{\mathbf m}_{g,\epsilon}^{(p)}
  =
  \relu\!\left(
    \mathbf U_{\eta_\epsilon}^{(p)}
    \mathbf z_{g,\rho'}^{(p-1)}
    +\mathbf b_{\graphrel,\eta_\epsilon}^{(p)}
  \right)
  \in\RR^d.
\label{eq:method:relation-message}
\end{equation}
The type-specific parameters satisfy
\begin{equation}
  \mathbf U_\eta^{(p)}\in\RR^{d\times d},
  \qquad
  \mathbf b_{\graphrel,\eta}^{(p)}\in\RR^d,
  \qquad
  \eta\in\mathcal T_{\graphrel}.
\label{eq:method:relation-edge-parameter-domains}
\end{equation}
We use \(\mathbf U\) for relation-graph matrices and \(\mathbf W\) for
row-graph matrices so that parameters from the two networks cannot be
confused. Every relation-edge weight equals one, and relation messages do
not contain an additional condition term.

For relation node \(\rho\), define its incoming edge-index set as
\begin{equation}
  \mathcal E_{g,\graphrel}^{-}(\rho)
  =
  \{\epsilon\in\mathcal E_g^{\graphrel}:\rgdst(\epsilon)=\rho\}.
\label{eq:method:relation-incoming-edges}
\end{equation}
The incoming messages are averaged using a denominator clamped at one:
\begin{equation}
  \widetilde{\mathbf a}_{g,\rho}^{(p)}
  =
  \frac{
    \displaystyle
    \sum_{\epsilon\in\mathcal E_{g,\graphrel}^{-}(\rho)}
    \widetilde{\mathbf m}_{g,\epsilon}^{(p)}
  }{
    \displaystyle
    \max\!\left(1,|\mathcal E_{g,\graphrel}^{-}(\rho)|\right)
  }
  \in\RR^d.
\label{eq:method:relation-aggregation}
\end{equation}
If \(\rho\) has no incoming relation edge, the numerator is the empty sum
\(\mathbf0\), the denominator is one, and hence the aggregate is
\(\mathbf0\).

The relation-node update combines the previous representation, a learned
self transformation, and the incoming aggregate:
\begin{equation}
  \mathbf z_{g,\rho}^{(p)}
  =
  \layernorm_{\graphrel}^{(p)}
  \left(
    \mathbf z_{g,\rho}^{(p-1)}
    +\relu\!\left(
      \mathbf U_{\mathrm{self}}^{(p)}
      \mathbf z_{g,\rho}^{(p-1)}
      +\mathbf b_{\graphrel,\mathrm{self}}^{(p)}
    \right)
    +\widetilde{\mathbf a}_{g,\rho}^{(p)}
  \right)
  \in\RR^d,
\label{eq:method:relation-update}
\end{equation}
where
\begin{equation}
  \mathbf U_{\mathrm{self}}^{(p)}\in\RR^{d\times d},
  \qquad
  \mathbf b_{\graphrel,\mathrm{self}}^{(p)}\in\RR^d.
\label{eq:method:relation-self-parameter-domains}
\end{equation}
For \(\boldsymbol\zeta\in\RR^d\), the layer-specific normalization is
\begin{equation}
  \layernorm_{\graphrel}^{(p)}(\boldsymbol\zeta)
  =
  \boldsymbol\gamma_{\graphrel}^{(p)}\odot
  \frac{
    \boldsymbol\zeta-\mu(\boldsymbol\zeta)\mathbf1
  }{
    \sqrt{\operatorname{var}(\boldsymbol\zeta)+\varepsilon_{\mathrm{ln}}}
  }
  +\boldsymbol\beta_{\graphrel}^{(p)},
\label{eq:method:relation-layernorm}
\end{equation}
with learned
\(\boldsymbol\gamma_{\graphrel}^{(p)},
\boldsymbol\beta_{\graphrel}^{(p)}\in\RR^d\) and
\(\varepsilon_{\mathrm{ln}}=10^{-5}\). Here
\(\mu(\boldsymbol\zeta)\) and
\(\operatorname{var}(\boldsymbol\zeta)\) are the mean and variance across
the \(d\) feature coordinates used by LayerNorm, and
\(\mathbf1\in\RR^d\) is the
all-ones vector. The four relation-edge types use
four different affine transformations within a layer, and the \(P=2\)
relation layers do not share their edge, self, or LayerNorm parameters.

\subsubsection{Relation-Conditioned Row-Graph Propagation}
\label{sec:method:row-propagation}

After the \(P=2\) relation layers, the final state of the relation node
associated with each retained FK edge is projected into the row-graph
hidden space. For
\(e\in\mathcal E_g^{\mathsf{f2p}}\cup
\mathcal E_g^{\mathsf{p2f}}\), define
\begin{equation}
  \boldsymbol\delta_e
  =
  \mathbf W_{\mathrm{cond}}
  \mathbf z_{g,\chi_g(e)}^{(P)}
  +\mathbf b_{\mathrm{cond}}
  \in\RR^d,
\label{eq:method:fk-condition}
\end{equation}
where
\begin{equation}
  \mathbf W_{\mathrm{cond}}\in\RR^{d\times d},
  \qquad
  \mathbf b_{\mathrm{cond}}\in\RR^d.
\label{eq:method:condition-parameter-domains}
\end{equation}
All row-level FK edges mapped to the same query-local relation node receive
the same relation state. In particular, the two directed copies of one FK
link satisfy
\begin{equation}
  \boldsymbol\delta_{\bar e}=\boldsymbol\delta_e.
\label{eq:method:shared-forward-reverse-condition}
\end{equation}
They nevertheless use different source-node states and different
\(\mathsf{f2p}\) and \(\mathsf{p2f}\) transformations. For every observed
label edge, we define \(\boldsymbol\delta_e=\mathbf0\). Both
\(\mathbf W_{\mathrm{cond}}\) and \(\mathbf b_{\mathrm{cond}}\) are
zero-initialized. The condition vectors are computed once after relation
propagation and remain fixed while the six row-graph layers update their
node states.

Under any fixed ordering in which row nodes precede candidate nodes, the
initial row-graph state matrix is
\begin{equation}
  \mathbf H_g^{(0)}
  =
  \operatorname{Concat}_{\mathrm{node}}
  \left(
    \{\mathbf h_{g,r}^{(0)}:r\in\mathcal R_g\},
    \{\mathbf h_{g,k}^{\mathrm{cand},(0)}:k\in\{0,1\}\}
  \right)
  \in\RR^{(|\mathcal R_g|+2)\times d}.
\label{eq:method:main-initial-state}
\end{equation}
Thus, candidate nodes are present during all message-passing layers rather
than being appended only at readout time.

Consider a directed row-graph edge
\(e:j\to i\), meaning
\(j=\rgsrc(e)\) and \(i=\rgdst(e)\). At layer
\(\ell\in\{1,\ldots,L\}\), its weighted, typed message is
\begin{equation}
  \mathbf m_{g,e}^{(\ell)}
  =
  \alpha_e\,
  \relu\!\left(
    \mathbf W_{\tau_e}^{(\ell)}
    \mathbf h_{g,j}^{(\ell-1)}
    +\mathbf b_{\tau_e}^{(\ell)}
    +\boldsymbol\delta_e
  \right)
  \in\RR^d,
\label{eq:method:main-message}
\end{equation}
where
\begin{equation}
  \mathbf W_\tau^{(\ell)}\in\RR^{d\times d},
  \qquad
  \mathbf b_\tau^{(\ell)}\in\RR^d,
  \qquad
  \tau\in\mathcal T_{\graphmain}.
\label{eq:method:main-edge-parameter-domains}
\end{equation}
The condition is added \emph{after} the type-specific affine transformation.
In particular, the implemented and non-implemented alternatives are
\begin{equation}
  \underbrace{
    \mathbf W_{\tau_e}^{(\ell)}\mathbf h_{g,j}^{(\ell-1)}
    +\mathbf b_{\tau_e}^{(\ell)}+\boldsymbol\delta_e
  }_{\text{implemented}}
  \ne
  \underbrace{
    \mathbf W_{\tau_e}^{(\ell)}
    (\mathbf h_{g,j}^{(\ell-1)}+\boldsymbol\delta_e)
    +\mathbf b_{\tau_e}^{(\ell)}
  }_{\text{not used}}.
\label{eq:method:condition-placement}
\end{equation}
The scalar \(\alpha_e\) multiplies the complete post-ReLU vector and is not
a component of the learned matrix \(\mathbf W_{\tau_e}^{(\ell)}\).

For node \(i\), define its incoming edge-index set and total incoming
weight as
\begin{equation}
\begin{aligned}
  \mathcal E_{g,\graphmain}^{-}(i)
  &=
  \{e\in\mathcal E_g^{\graphmain}:\rgdst(e)=i\},\\
  \deg_{g,\alpha}(i)
  &=
  \sum_{e\in\mathcal E_{g,\graphmain}^{-}(i)}\alpha_e.
\end{aligned}
\label{eq:method:main-incoming-degree}
\end{equation}
Messages from \(\mathsf{f2p}\), \(\mathsf{p2f}\), and
\(\mathsf{label}\) edges enter a single shared aggregate:
\begin{equation}
  \mathbf a_{g,i}^{(\ell)}
  =
  \frac{
    \displaystyle
    \sum_{e\in\mathcal E_{g,\graphmain}^{-}(i)}
    \mathbf m_{g,e}^{(\ell)}
  }{
    \displaystyle
    \max\!\left(1,\deg_{g,\alpha}(i)\right)
  }
  \in\RR^d.
\label{eq:method:main-aggregation}
\end{equation}
The denominator is the combined weight of all three incoming edge types,
not a separate normalization for each type. The lower bound avoids division
by zero and means that the expression is not a strict weighted average when
\(0<\deg_{g,\alpha}(i)<1\). If there is no incoming edge, the empty numerator
and clamped denominator give \(\mathbf a_{g,i}^{(\ell)}=\mathbf0\).

The node update consists of an identity residual, a learned self
transformation, and the incoming aggregate:
\begin{equation}
  \mathbf h_{g,i}^{(\ell)}
  =
  \layernorm_{\graphmain}^{(\ell)}
  \left(
    \mathbf h_{g,i}^{(\ell-1)}
    +\relu\!\left(
      \mathbf W_{\mathrm{self}}^{(\ell)}
      \mathbf h_{g,i}^{(\ell-1)}
      +\mathbf b_{\mathrm{self}}^{(\ell)}
    \right)
    +\mathbf a_{g,i}^{(\ell)}
  \right)
  \in\RR^d,
\label{eq:method:main-update}
\end{equation}
where
\begin{equation}
  \mathbf W_{\mathrm{self}}^{(\ell)}\in\RR^{d\times d},
  \qquad
  \mathbf b_{\mathrm{self}}^{(\ell)}\in\RR^d.
\label{eq:method:main-self-parameter-domains}
\end{equation}
The main LayerNorm is defined analogously to
Eq.~\eqref{eq:method:relation-layernorm}, with independent learned
\(\boldsymbol\gamma_{\graphmain}^{(\ell)},
\boldsymbol\beta_{\graphmain}^{(\ell)}\in\RR^d\) for each layer.

The graph structure, scalar weights, and conditions stay fixed across
layers, while messages are recomputed from the newly updated source states:
\begin{equation}
  \mathbf H_g^{(\ell)}
  =
  \operatorname{MainMPLayer}^{(\ell)}
  \left(
    \mathbf H_g^{(\ell-1)},
    \mathcal G_g^{\graphmain},
    \{\boldsymbol\delta_e\}_{e\in\mathcal E_g^{\graphmain}}
  \right),
  \qquad \ell=1,\ldots,L.
\label{eq:method:main-layer-stack}
\end{equation}
For \(L=6\), every layer owns a distinct self transform, three distinct
edge-type transforms, and its own LayerNorm. Candidate nodes have no FK
neighbors and no outgoing edges. Their incoming aggregate can therefore
contain only observed-label messages; when no such edge is present, they
still evolve through the residual and self-transformation paths. There is
no candidate-to-row feedback.

The computation order is strictly
\begin{equation}
  \underbrace{\mathcal G_g^{\graphrel}
  \xrightarrow{\ P=2\ \text{layers}\ }
  \{\boldsymbol\delta_e\}}_{\text{relation encoding}}
  \quad\longrightarrow\quad
  \underbrace{\mathcal G_g^{\graphmain}
  \xrightarrow{\ L=6\ \text{layers}\ }
  \mathbf H_g^{(L)}}_{\text{row and candidate encoding}}.
\label{eq:method:two-stage-propagation}
\end{equation}
The relation network is not interleaved between row-graph layers.

\subsubsection{Base Model as a Special Case}
\label{sec:method:base-special-case}

The base and relation-conditioned variants share the sampled context, main
graph, cell, row, and candidate initialization, six-layer row GNN, readout, and
classification loss. Their only difference inside the row GNN is whether an
FK message receives \(\boldsymbol\delta_e\). For an FK edge, the base model
computes the following condition-free message:
\begin{equation}
  \mathbf m_{g,e,\mathrm{base}}^{(\ell)}
  =
  \alpha_e\,
  \relu\!\left(
    \mathbf W_{\tau_e}^{(\ell)}
    \mathbf h_{g,\rgsrc(e)}^{(\ell-1)}
    +\mathbf b_{\tau_e}^{(\ell)}
  \right).
\label{eq:method:base-message}
\end{equation}
In the relation-conditioned model, an FK edge instead uses
\begin{equation}
  \mathbf m_{g,e,\mathrm{rel}}^{(\ell)}
  =
  \alpha_e\,
  \relu\!\left(
    \mathbf W_{\tau_e}^{(\ell)}
    \mathbf h_{g,\rgsrc(e)}^{(\ell-1)}
    +\mathbf b_{\tau_e}^{(\ell)}
    +\boldsymbol\delta_e
  \right).
\label{eq:method:rel-message}
\end{equation}
Equivalently, both cases are summarized by
\begin{equation}
  \boldsymbol\delta_e
  =
  \begin{cases}
    \mathbf0,
      & \text{base model or observed-label edge},\\
    \mathbf W_{\mathrm{cond}}
      \mathbf z_{g,\chi_g(e)}^{(P)}+\mathbf b_{\mathrm{cond}},
      & \text{relation-conditioned FK edge}.
  \end{cases}
\label{eq:method:base-rel-condition}
\end{equation}
At code level, disabling the relation graph omits its query marker,
relation layers, and condition projection; the main MPLayer is then called
without an FK-condition tensor. Label edges are unconditioned in both
variants. Because the condition projection is zero-initialized, a
relation-conditioned model initially has zero FK conditions and matches the
base message function provided the two models' main-graph parameters are
the same. If a sampled query contains no FK edge, the auxiliary relation
branch has no effect.

\subsection{Candidate Scoring and Classification Objective}
\label{sec:method:classification-objective}

After \(L=6\) row-graph layers, the final seed representation is
\(\mathbf h_{g,r_g^\star}^{(L)}\in\RR^d\), and the final representation of
candidate \(k\in\{0,1\}\) is
\(\mathbf h_{g,k}^{\mathrm{cand},(L)}\in\RR^d\). For each candidate, the
readout input concatenates the query context with the output hypothesis:
\begin{equation}
  \boldsymbol\xi_{g,k}^{\mathrm{ro}}
  =
  [
    \mathbf h_{g,r_g^\star}^{(L)}
    \vcat
    \mathbf h_{g,k}^{\mathrm{cand},(L)}
  ]
  \in\RR^{2d}.
\label{eq:method:readout-input}
\end{equation}
A two-layer MLP first produces the hidden readout representation
\begin{equation}
  \boldsymbol\psi_{g,k}^{\mathrm{ro}}
  =
  \relu\!\left(
    \mathbf W_{\mathrm{ro},1}
    \boldsymbol\xi_{g,k}^{\mathrm{ro}}
    +\mathbf b_{\mathrm{ro},1}
  \right)
  \in\RR^d,
\label{eq:method:readout-hidden}
\end{equation}
and then a scalar candidate score
\begin{equation}
  s_{g,k}
  =
  \mathbf w_{\mathrm{ro},2}^{\top}
  \boldsymbol\psi_{g,k}^{\mathrm{ro}}
  +b_{\mathrm{ro},2}
  \in\RR.
\label{eq:method:candidate-score}
\end{equation}
The readout parameters have domains
\begin{equation}
  \mathbf W_{\mathrm{ro},1}\in\RR^{d\times2d},
  \qquad
  \mathbf b_{\mathrm{ro},1}\in\RR^d,
  \qquad
  \mathbf w_{\mathrm{ro},2}\in\RR^d,
  \qquad
  b_{\mathrm{ro},2}\in\RR.
\label{eq:method:readout-parameter-domains}
\end{equation}
Exactly the same MLP parameters score False and True candidates for every
task and query; the model has no task-specific prediction head. Thus,
\(s_{g,k}\) is a raw seed--candidate compatibility score rather than a
probability.

For binary classification, define the single logit as the difference
between the True and False scores:
\begin{equation}
  \Delta_g
  =
  s_{g,1}-s_{g,0}
  \in\RR.
\label{eq:method:classification-logit}
\end{equation}
The probability assigned to the True class is
\begin{equation}
  \pi_g
  =
  \operatorname{sigmoid}(\Delta_g)
  =
  \frac{\exp(s_{g,1})}{\exp(s_{g,0})+\exp(s_{g,1})}
  =
  \Pr(y_g=1\mid\mathcal X_g),
\label{eq:method:classification-probability}
\end{equation}
where \(\mathcal X_g\) denotes the complete constructed input for query
\(g\), including the row graph and, for the relation-conditioned variant,
its FK conditions. The False probability is \(1-\pi_g\). The corresponding
hard prediction selects the higher-scoring candidate:
\begin{equation}
  \widehat y_g
  =
  \underset{k\in\{0,1\}}{\arg\max}\;s_{g,k}.
\label{eq:method:classification-prediction}
\end{equation}
Therefore, \(s_{g,1}>s_{g,0}\) predicts True, whereas
\(s_{g,0}>s_{g,1}\) predicts False. Ranking-based evaluation can operate
directly on \(\Delta_g\) without first thresholding it.

The per-query binary cross-entropy is
\begin{equation}
  \mathcal L_g^{\mathrm{clf}}
  =
  -y_g\log\pi_g
  -(1-y_g)\log(1-\pi_g).
\label{eq:method:classification-loss}
\end{equation}
The implementation evaluates this expression with a numerically stable
BCE-with-logits routine applied directly to \(\Delta_g\). Let \(\mathcal B\)
be the set of valid binary query graphs in a classification-only minibatch,
and define \(G=|\mathcal B|\). Their mean loss is
\begin{equation}
  \mathcal L_{\mathcal B}^{\mathrm{clf}}
  =
  \frac{1}{G}
  \sum_{g\in\mathcal B}\mathcal L_g^{\mathrm{clf}}.
\label{eq:method:classification-batch-loss}
\end{equation}
Every valid query graph has equal weight, irrespective of its number of
sampled cells, row nodes, or edges. Because the scalar output bias is shared
between candidates, it cancels in the classification logit:
\(b_{\mathrm{ro},2}-b_{\mathrm{ro},2}=0\). It is retained in the score
definition to match the shared readout implementation and to support the
future multi-candidate regression extension.

\subsection{Extension to Regression}
\label{sec:method:regression-placeholder}

\emph{Placeholder. This subsection will define continuous candidate
construction, observed-label-to-candidate edges, the regression objective,
continuous inference, and the mixed classification--regression minibatch
loss after the regression design is finalized.}
